\documentclass[10pt,a4paper]{amsart}
\usepackage[margin=19mm]{geometry}
\usepackage{amsmath,amssymb,mathtools}
\usepackage[scr=rsfs,cal=euler]{mathalfa}
\usepackage{xfrac}
\usepackage[unicode,hidelinks,bookmarks=false]{hyperref}
\newcommand{\ie}{i.e.\ }
\newcommand{\cf}{cf.\ }
\newcommand{\resp}{respectively\ }
\newcommand{\Subsection}{\subsection}
\providecommand{\nobreakdash}{\nobreak\hbox{-}}
\setlength{\emergencystretch}{2em}
\multlinegap0pt
\usepackage{amssymb}
\usepackage{mathtools}
\usepackage{float}
\usepackage[shortlabels]{enumitem}
\usepackage{xcolor}
\usepackage{bm}
\newtheorem{lemma}{Lemma}
\newcommand{\bR}{\mathbb{R}}
\newcommand{\cI}{\mathcal{I}}
\newcommand{\cl}{\colon}
\title{Open-closed Deligne--Mumford field theories: geometric foundations}
\author{Amanda Hirschi, Kai Hugtenburg}
\date{Source: arXiv:2501.04687v2 (CC BY 4.0)}
\begin{document}
\maketitle

\noindent\emph{Source and licence.} Open-closed Deligne--Mumford field theories: geometric foundations. Version arXiv:2501.04687v2 (CC BY 4.0), distributed by arXiv under the Creative Commons Attribution 4.0 International licence, http://creativecommons.org/licenses/by/4.0/, as recorded in the arXiv licence metadata for this exact version and shown on the abstract page https://arxiv.org/abs/2501.04687v2. Full licence notice: \texttt{licenses/arxiv-2501.04687-CC-BY-4.0.txt}.\par\smallskip

\noindent\textit{Editorial scope.} Counted unit: the article's lemma delta (source0.tex lines 2953--2956). The article's complete original proof of it is delivered with it (source0.tex lines 2958--3006, one \texttt{proof} environment of 49 lines). No mathematics is written, repaired or re-derived: the statement and the proof are the article's own lines, in the article's own notation, and no retained line is substituted or re-flowed.  No further source-local context is needed: every cross-reference of every retained line resolves inside the two delivered spans.  Nothing was omitted from any retained span, so the package carries no omission record.  No retained line contains a citation, so the wrapper carries no bibliography and no bibliography entry.  No retained line contains a figure environment, an image-inclusion command or drawing code, so no image is delivered and the package declares no asset dependency.  Editorial formatting only, in addition: an AMS \texttt{amsart} preamble that restates the article's own declarations and macros used by the retained lines, namely the environments \texttt{lemma} on separate counters, together with the standard packages those lines need.  The retained environments print in this document as Lemma 1 in order of appearance, whereas the article prints the same items with the numbering of its own full document.  The complete native source of the article as distributed in the arXiv e-print for this exact version is bundled unchanged as \texttt{sources/171354/source0.tex}.
\begin{lemma}
	\label{lem:local exactness of delta}
	Suppose $M = \bR^n \times Y$ for some smooth manifold $Y$ and that $H_i \subset \bR^n$ is a linear subspace (not necessarily of codimension 1) for each $i\in\cI$. Assume $M_i = H_i \times Y$ and let $E \rightarrow Y$ be any vector bundle. Then the differential $\delta$ is exact on $C^{*,*}(\{H_i \times Y\}, \pi^*E)$, where $\pi\cl  \bR^n \times Y \rightarrow Y$ is the projection.
\end{lemma}

\begin{proof}
	Let $f \in C^{*,q}(\{H_i \times Y\}, \pi^*E)$ be $\delta$-closed. Then for $J = (j_1, \dots, j_q)$, define
	\begin{equation*}
		g_{J} = \sum_{I<k<j_1} (-1)^{|I|}\pi_{J}^{IkJ} \circ i_{IkJ} f_{kJ},
	\end{equation*}
	where $\pi_{J}^{IkJ}\cl \Omega^*_{cv}(E|_{M_{IkJ}}) \rightarrow \Omega^*_{cv}(E|_{M_{J}})$ is induced by the projection map $M_{J} \rightarrow M_{IkJ}$, and the sum is over all $k \notin J$ and $I = (i_1, \dots, i_r) \subset \{1, \dots n\} \setminus (J \cup \{k\})$ such that the tuple $(i_1, \dots, i_r, k, j_1, \dots, j_q)$ is ordered. We allow for $I =  \emptyset$. Then, for any $J = (j_0, \dots, j_q)$, we have 
	\begin{align}	\label{eq:terms with p>0}
		\notag\delta(g)_{J} &:= \sum_p (-1)^p i_{J} g_{\hat{J}_p}\\
		&= \sum_{\substack{p \neq 0\\I < k < j_0}} (-1)^{p+|I|} \pi_{J}^{IkJ} \circ i_{IkJ} f_{k\hat{J}_p}+ \sum_{I < k < j_1} (-1)^{|I|} \pi_{J}^{IkJ} \circ i_{IkJ} f_{k\hat{J}_0},
	\end{align}
	where we have separated the contributions from $p=0$ and $p >0$, and used the fact that 
	$$i_{J} \circ \pi_{\hat{J}_p}^{Ik \hat{J}_p} \circ i_{Ik \hat{J}_p} = \pi_{J}^{IkJ} \circ i_{Ik J}.$$ 
	We next separate the contributions with $p=0$ further by taking terms 
	\begin{itemize}
		\item with $(k,I) = (j_0, \emptyset)$,
		\item with $k = j_0$ and $|I| >1$
		\item with $k < j_0$,
		\item and with $k>j_0$,
	\end{itemize} 
	obtaining
	\begin{align}
		\sum_{I < k < j_1} (-1)^{|I|} \pi_{J}^{IkJ} \circ i_{IkJ} f_{j \hat{J}_0} = f_J
		&+\sum_{\substack{I\neq \emptyset\\I < j_0}} (-1)^{|I|} \pi_{J}^{IJ} \circ i_{IJ} f_{J} 
		\label{eq:terms with I nonempty and k = j0}\\
		&+\sum_{I < k < j_0} (-1)^{|I|} \pi_{J}^{IkJ} \circ i_{IkJ} f_{j \hat{J}_0}
		\label{eq:terms with k < j0}\\
		&+\sum_{\substack{k > j_0, I\\(I,k,j_1)}} (-1)^{|I|} \pi_{J}^{IkJ} \circ i_{IkJ}f_{k\hat{J}_0}.
		\label{eq:terms with k > j0}
	\end{align}
	First we show that the final contribution in line \eqref{eq:terms with k > j0} vanishes, computing
	\begin{align*}
		\sum_{\substack{k > j_0, I\\I < k < j_1}} (-1)^{|I|} \pi_{J}^{IkJ} \circ i_{IkJ} f_{k\hat{J}_0} &= \sum_{\substack{k > j_0, I\\I < k < j_1 \\ 
				j_0 \in I}} (-1)^{|I|} \pi_{J}^{IkJ} \circ i_{IkJ} f_{k\hat{J}_0} + \sum_{\substack{k > j_0, I\\I < k < j_1 \\ j_0 \notin I}} (-1)^{|I|} \pi_{J}^{IkJ} \circ i_{IkJ}f_{k\hat{J}_0}\\
		&= \sum_{\substack{k > j_0, I\\I < k < j_0\\ j_0 \notin I}} (-1)^{|I| + 1} \pi_{J}^{IkJ} \circ i_{IkJ} f_{k\hat{J}_0} + \sum_{\substack{k > j_0, I\\I < k< j_1\\ j_0 \notin I}} (-1)^{|I|} \pi_{J}^{IkJ} \circ i_{IkJ}f_{k\hat{J}_0}\\& = 0.
	\end{align*}
	The second equality holds as $\pi_{J}^{IkJ} \circ i_{IkJ}$ does not depend on whether $j_0 \in I$ since in either case we set the variable $x_{j_0} = 0$. Next, we note that 
	\begin{align*}
		0 = \delta(f)_{kJ} = i_{kJ} f_J + \sum_{p = 0}^k (-1)^{p+1} i_{kJ} f_{k\hat{J}_p}.
	\end{align*}
	Applying $\pi_{J}^{IkJ} i_{IkJ}$ to this equation and separating $p = 0$ and $p > 0$ yields
	\begin{equation*}
		0 = \pi_{J}^{IkJ} i_{IkJ} f_J - \pi_{J}^{IkJ} i_{IkJ} f_{k\hat{J}_0} + \sum_{p=1}^{k} (-1)^{p+1} \pi_{J}^{IkJ} i_{IkJ} f_{k\hat{J}_p}.
	\end{equation*}
	Next we rewrite the terms with $I \neq \emptyset$ and $k = j_0$ in line \eqref{eq:terms with I nonempty and k = j0}, that is,
	\begin{equation*}
		\sum_{\substack{I\neq \emptyset\\I < j_0}} (-1)^{|I|} \pi_{J}^{IJ} \circ i_{IJ} f_{J} = \sum_{I < k < j_0} (-1)^{|I|+1} \pi_{J}^{IkJ} \circ i_{IkJ} f_{J}.
	\end{equation*}
	We thus find that the terms with $p > 0$ in line \eqref{eq:terms with p>0}, those with $p = 0, I \neq \emptyset, k = j_0$ in line \eqref{eq:terms with I nonempty and k = j0}, and those with $p =0, k < j_0$ in line \eqref{eq:terms with k < j0} are exactly covered by the relation $\delta(f) = 0$.
\end{proof}

\end{document}
